% SEFM 2026 Tool Paper - Springer LNCS format
% Temporarily using article class for compilation
% Final submission will use: \documentclass[runningheads]{llncs}
\documentclass[11pt,a4paper]{article}
\usepackage[margin=1in]{geometry}

\usepackage{graphicx}
\usepackage{listings}
\usepackage{hyperref}
\usepackage{xspace}

% Tool name command
\newcommand{\tool}{\texttt{rocq-piler}\xspace}
\newcommand{\axm}{\texttt{axiomander}\xspace}

\begin{document}

\title{\tool: A Content-Addressed Proof Oracle for\\ LLM-Driven Discharge of Separation-Logic Obligations}

\author{Gavin Mendel-Gleason\\
Scidonia, Dublin, Ireland\\
\texttt{gavin@scidonia.com}}

\maketitle

\begin{abstract}
We present \tool, a tool that connects Large Language Models (LLMs) to the Rocq
proof assistant (formerly Coq) through the Model Context Protocol (MCP), and that
is designed around a single organising idea: \emph{content-addressed proof
obligations}. Every open goal in a proof is identified by a stable hash of its
statement and hypotheses, so an LLM can target, batch, and re-target obligations
by hash rather than by fragile line numbers or bullet positions. On top of this
substrate \tool provides higher-level moves---\texttt{stratify} (case-split a
proof and discharge the easy cases with a tactic portfolio) and
\texttt{close\_admits} (batch-close the survivors)---together with speculative
execution and an incremental verification cache. Our motivating application is
the deductive verification system \axm, which compiles imperative programs to
Iris separation-logic proof obligations. Automation discharges most of these
obligations, but a residual fraction (in our experience roughly one in five)
requires an expert. \tool is the \emph{oracle} we are tuning to play that
expert role. As a self-contained, reproducible case study we report an
autonomous proof of type preservation for PCF with mutable references (21
inductive cases, 9 supporting lemmas), completed by two different LLMs through
\tool's MCP interface at a cost of a few cents per run. The tool is open-source
and under active development.
\end{abstract}

\noindent\textbf{Keywords}: Proof assistants, Rocq, Coq, separation logic, Iris,
LLMs, deductive verification, proof automation, MCP, neurosymbolic AI

\section{Introduction}

Deductive program verifiers such as those built on Iris~\cite{jung2018iris}
reduce the correctness of imperative programs to a set of \emph{proof
obligations} in separation logic. Powerful automation can dispatch many of these
obligations, but not all: a residual tail of obligations---arising from
non-trivial framing, ghost state, invariants, or arithmetic side conditions---is
left for a human expert. This tail is precisely where verification projects stall,
because it demands scarce expertise and tedious interactive proof.

\tool targets this tail. It is a tool that exposes the Rocq proof
assistant~\cite{coq} to Large Language Models (LLMs) through the Model Context
Protocol (MCP), with the explicit goal of acting as an \emph{oracle} that an
automated verifier can call when its own tactics give up. The LLM proposes proof
steps based on learned patterns; Rocq, accessed through its Language Server
Protocol interface~\cite{gallego2022coqlsp} and the Pétanque execution backend,
verifies every step. This division of labour---neural proposal, symbolic
verification---means the LLM need never be trusted: unsound steps are simply
rejected.

The novel element of \tool is not the bridge itself but its \emph{interface
abstraction}. Earlier LLM--prover integrations address goals positionally
(line/column, bullet nesting), which is brittle: a single inserted tactic
invalidates every downstream coordinate, and multi-goal proofs quickly become
unmanageable for an agent. \tool instead makes every open obligation
\emph{content-addressed}: each goal is identified by a stable hash derived from
its statement and hypotheses. Identical goals share a hash and can be closed
together; an obligation can be revisited later by the same hash even after the
surrounding script has changed. This turns proof construction into a simple
loop---\emph{enumerate hashes, write a tactic for each hash, repeat until
\texttt{Qed}}---that is robust under edits and natural for an autonomous agent.

\paragraph{Contributions.}
\begin{itemize}
\item A \textbf{content-addressed obligation model} for interactive proof, in
  which open goals are referenced by stable hashes rather than positions
  (Section~\ref{sec:hashes}).
\item Two \textbf{batch proof moves} built on this model---\texttt{stratify}
  (skeleton case-split plus a per-subgoal closing portfolio) and
  \texttt{close\_admits} (portfolio-based batch closing of survivors)---plus
  speculative execution and an \textbf{incremental verification cache} that
  re-checks only changed functions and their affected callers
  (Section~\ref{sec:tools}).
\item An account of \tool as a \textbf{proof oracle for \axm}, an open
  deductive verifier that emits Iris separation-logic obligations, where roughly
  one fifth of obligations resist automation and fall to the oracle
  (Section~\ref{sec:axiomander}).
\item A \textbf{reproducible case study}: autonomous proof of type preservation
  for PCF with references (21 cases, 9 lemmas) by two distinct LLMs, with cost
  and tool-call statistics (Section~\ref{sec:casestudy}).
\end{itemize}

\textbf{Availability}: \tool is open-source (MIT license), available at
\url{https://github.com/scidonia/rocq-piler} and published on npm as
\texttt{rocq-piler}.

\section{Motivation: the Residual Obligations of a Deductive Verifier}
\label{sec:axiomander}

\axm is a separate open-source deductive verification system. It takes annotated
imperative code and produces proof obligations in Iris, a higher-order
concurrent separation logic embedded in Rocq~\cite{jung2018iris}. Like other
verifiers in this family, \axm discharges a large majority of its obligations
automatically through built-in tactics and heuristics. The difficulty is the
remainder.

In our experience, automation alone leaves on the order of \textbf{20\% of
obligations} unproved. These are the obligations that require genuine proof
engineering: choosing how to frame a heap assertion, supplying an existential
witness for ghost state, strengthening a loop invariant, or closing an
arithmetic side goal that the automation cannot see through. Historically this
20\% is handled by a human expert working interactively in Rocq---the bottleneck
that limits how much code a verification effort can cover.

\tool is our attempt to automate that expert. Rather than extending \axm's
internal automation, we treat the leftover obligations as ordinary Rocq goals
and hand them to an LLM-driven oracle that interacts with Rocq through \tool. The
verifier remains responsible for soundness---every obligation is still a Rocq
goal checked by Rocq's kernel---while the oracle supplies the creative,
hard-to-automate proof steps. Because the oracle is an LLM behind an MCP
interface, it can be \emph{tuned}: the tool surface, the prompting, the closing
portfolios, and the model itself are all parameters we adjust to raise the
fraction of the residual obligations that close without a human.

This integration is an early-stage prototype. In this paper we therefore use a
self-contained Rocq development---type preservation for PCF with
references---as a faithful, reproducible stand-in for the kind of obligation the
oracle must close: deep induction, auxiliary lemmas about substitution and
weakening, and reasoning about a mutable store. It exercises exactly the
oracle behaviours that the \axm tail requires, without depending on the full
\axm pipeline.

\section{The Content-Addressed Obligation Model}
\label{sec:hashes}

\subsection{From positions to hashes}

A Rocq proof in progress is a tree of open goals beneath a partially written
script. Addressing those goals by position is fragile: inserting a tactic shifts
every subsequent line, and Rocq's bullet discipline ($-$, $+$, $*$, braces)
makes the ``current'' goal depend on subtle structural context. For an
autonomous agent that edits the script repeatedly, positional addressing is a
constant source of stale references.

\tool replaces positions with content. Each open obligation is sealed as a named
\texttt{admit} annotated with an 8-character hash computed from the goal's
statement and its hypothesis context. The hash has two properties that drive the
whole workflow:

\begin{itemize}
\item \textbf{Stability}: the hash depends only on the goal, not on where it sits
  in the script. An obligation can be revisited by the same hash after unrelated
  edits elsewhere.
\item \textbf{Coalescence}: syntactically identical goals share a hash. A proof
  that produces four copies of the goal \texttt{True}, possibly at different
  bullet depths, exposes a single hash that closes all four at once.
\end{itemize}

\subsection{The core loop}

Proof construction with \tool reduces to a short cycle:

\begin{enumerate}
\item \texttt{focus\_proof} returns the proof state, the script so far, and the
  list of open obligations---each with its hash, hypotheses, and goal.
\item For each hash, \texttt{insert\_tactics} commits a tactic (or a multi-line
  block) targeting that obligation by hash. If the tactic leaves several goals,
  they are re-sealed as fresh hash-addressable admits.
\item Repeat until no obligations remain, at which point the proof is closed
  automatically with \texttt{Qed}.
\end{enumerate}

Crucially, the agent reads hypotheses-plus-goal for every obligation directly
from the \texttt{focus\_proof} (or \texttt{insert\_tactics}) response, so it can
write each case's tactic without re-querying the prover for context.

\section{Tool Architecture and Higher-Level Moves}
\label{sec:tools}

\subsection{System overview}

\tool is implemented in TypeScript on top of three layers:

\begin{enumerate}
\item an \textbf{MCP server} exposing the tool surface to any MCP-compatible
  client (e.g.\ Claude, OpenCode, custom agents);
\item an \textbf{LSP/Pétanque client} that drives \texttt{rocq-lsp} for document
  synchronisation, diagnostics, and speculative state execution; and
\item a \textbf{workspace manager} that detects project roots
  (\texttt{\_CoqProject}, \texttt{\_RocqProject}, \texttt{dune-project}) and
  keeps the language server aligned with the right project.
\end{enumerate}

\subsection{Core obligation tools}

\begin{itemize}
\item \texttt{focus\_proof}: one-stop entry point---returns goal state, the
  script up to the cursor, and every open admit with its hash, hypotheses, and
  goal.
\item \texttt{insert\_tactics}: insert a tactic or multi-line block; target a
  specific obligation with \texttt{admit\_hash}; supports speculative
  \texttt{dry\_run} evaluation and \texttt{replace} for retry.
\item \texttt{open\_goals} / \texttt{proof\_state}: lighter-weight goal queries.
\item \texttt{add\_lemma} / \texttt{delete\_lemma}: insert or remove helper
  lemma stubs above a target proof.
\item \texttt{reset\_proof} / \texttt{undo\_step} / \texttt{delete\_step}:
  structured rollback of a proof body or of recent edits.
\end{itemize}

\subsection{Batch moves: \texttt{stratify} and \texttt{close\_admits}}

Two moves make multi-case proofs tractable for an agent.

\paragraph{\texttt{stratify}} runs a \emph{skeleton} tactic that case-splits the
proof (for example, \texttt{induction H; intros; inversion Ht; subst}), then
attempts each resulting subgoal independently with every tactic in a
\emph{closing portfolio}, each wrapped in \texttt{solve[\ldots]} so that partial
progress never leaks between cases. Solved cases have their winning tactic
inlined; unsolved cases become labelled, hash-addressable admit bullets. A
single \texttt{stratify} call can thus dispatch the routine majority of an
induction and report only the genuine survivors, by hash.

\paragraph{\texttt{close\_admits}} takes a portfolio mapping hashes (or the
wildcard \texttt{*}) to closing tactics and applies them speculatively: a tactic
is committed only if it closes its target obligation, otherwise the obligation is
left untouched. This lets the agent sweep up survivors in one call while
guaranteeing the script never regresses.

Together these moves implement the natural strategy for a hard goal: split it,
auto-close what you can, and spend the LLM's attention only on the residue---the
same shape as the \axm problem of automating the 20\% tail.

\subsection{Speculative execution and knowledge search}

\tool exposes the Pétanque API (\texttt{snap\_state}, \texttt{exec\_tactic},
\texttt{state\_goals}) and a \texttt{dry\_run} mode so an agent can test a tactic
without modifying the file. Knowledge tools---\texttt{search\_lemmas},
\texttt{inspect\_term}, \texttt{inspect\_about}, \texttt{locate\_term},
\texttt{require\_lib}---let the agent discover applicable lemmas and definitions
before committing.

\subsection{Incremental verification cache}

Because an oracle is invoked repeatedly over an evolving development, \tool
caches verification at the granularity of individual functions/proofs. It tracks
separate hashes for a proof's body, its contract/statement, and the contracts of
its callees, and re-verifies only what changed plus the transitive callers
affected by a contract change. Unchanged, previously verified obligations return
instantly. This keeps the cost of repeatedly consulting the oracle low---an
essential property when it is wired into a verifier's inner loop.

\section{Case Study: Type Preservation for PCF with References}
\label{sec:casestudy}

\subsection{Problem}

As a reproducible proxy for the \axm residual obligations, we verify type
preservation for PCF extended with mutable references:
\[
\frac{\Gamma; \Sigma \vdash e : \tau \quad (e,\mu) \rightarrow (e',\mu')
      \quad \mu : \Sigma}
     {\exists\,\Sigma' \supseteq \Sigma.\;\; \mu' : \Sigma'
      \;\wedge\; \Gamma; \Sigma' \vdash e' : \tau}
\]
The development includes a typed store, store-typing extension, and a de~Bruijn
substitution calculus. The proof requires 21 inductive step cases and 9
supporting lemmas covering store extension (\texttt{extends\_refl},
\texttt{extends\_trans}, \texttt{nth\_error\_extends}), weakening
(\texttt{store\_weakening}, \texttt{heap\_ok\_extends}), shifting and
substitution (\texttt{shift\_typing}, \texttt{subst\_preserves\_typing}), and
heap invariants (\texttt{heap\_lookup\_has\_type}, \texttt{heap\_update\_ok}).

\subsection{Autonomous process}

The proof was driven entirely through \tool's MCP interface. The agent
proceeded as the architecture intends:

\begin{enumerate}
\item Introduce auxiliary lemma stubs with \texttt{add\_lemma} and prove each by
  the focus/insert loop; the substitution and shifting lemmas required a
  generalised induction (over an arbitrary context prefix) that the agent
  discovered after an initial formulation failed and was reset.
\item Attack preservation with a single \texttt{stratify}: a skeleton induction
  over the step relation plus a closing portfolio, which discharged 18 of the
  21 cases automatically and left three survivors (\texttt{S\_AppAbs},
  \texttt{S\_DerefLoc}, \texttt{S\_AssignV}) addressed by hash.
\item Close the three survivors individually with \texttt{insert\_tactics}
  targeted by \texttt{admit\_hash}, invoking the substitution and heap lemmas.
\end{enumerate}

\subsection{Results}

The full development closes with \texttt{Qed} on all 9 lemmas and the main
theorem, with no remaining admits. We ran the task with two distinct
models---an open-weight model (DeepSeek~V4) and a proprietary model (a Claude
release)---each completing the proof through the same MCP interface. A
representative DeepSeek~V4 run used \textbf{50 API calls} at an actual provider
cost of \textbf{about four cents} (\$0.04). The \texttt{stratify}/\texttt{close\_admits}
pattern was decisive: the bulk of the 21 cases never reached the model as
individual problems, mirroring the goal of letting the oracle spend effort only
on the genuinely hard obligations.

\begin{table}[t]
\centering
\caption{Selected case-study statistics (DeepSeek~V4 run).}
\begin{tabular}{lr}
\hline
Inductive step cases & 21\\
Auxiliary lemmas & 9\\
Cases auto-closed by a single \texttt{stratify} & 18\\
Hash-addressed survivors closed individually & 3\\
Total API calls & 50\\
Actual provider cost & \$0.04\\
\hline
\end{tabular}
\end{table}

\section{Related Work}

\textbf{Separation-logic verification.} Iris~\cite{jung2018iris} underpins a
family of deductive verifiers whose automation leaves a residual tail of
interactive obligations; \axm sits in this family, and \tool targets that tail
rather than the logic itself.

\textbf{LLM-based theorem proving.} Baldur~\cite{first2023baldur} generates and
repairs Coq proofs with retrieval-augmented LLMs, and
LeanDojo~\cite{yang2023leandojo} provides a learning environment for Lean. These
focus on model training and bespoke environments; \tool instead standardises the
\emph{interaction protocol} (MCP) and contributes the content-addressed
obligation model and batch moves, so any MCP-capable model can serve as the
oracle.

\textbf{Prover interfaces.} \texttt{coq-lsp}~\cite{gallego2022coqlsp} provides
the LSP foundation \tool builds on; earlier IDE-oriented tools such as
CoqPIE~\cite{roe2016coqpie} and Proof General~\cite{aspinall2000proof} target
human users. \tool adds an agent-oriented layer: hash-addressed goals,
speculative execution, batch closing, and an incremental cache.

\section{Status, Lessons, and Future Work}

The \axm integration is a working prototype: leftover Iris obligations can be
routed to the oracle, and we are actively tuning the tool surface, prompting,
and closing portfolios to increase the share of the 20\% tail that closes
without a human. Two design lessons stand out. First, \emph{addressing
obligations by content rather than position} is what makes autonomous,
multi-goal proof construction robust to the agent's own edits. Second,
\emph{batch moves with speculative, non-regressing semantics}
(\texttt{stratify}, \texttt{close\_admits}) match how the residual-obligation
problem actually decomposes: auto-close the routine majority, escalate only the
survivors.

Future work includes (i) end-to-end evaluation on real \axm obligation suites
with success-rate and cost reporting; (ii) feeding verifier-side context (framing
hints, ghost-state signatures) into the oracle's prompts; and (iii) learning
closing portfolios from past successful discharges.

\section{Conclusion}

\tool reframes LLM-driven proving around content-addressed obligations and
batch, speculative closing moves, and positions the result as a tunable oracle
for the hard tail of a separation-logic verifier, \axm. A reproducible PCF+references
case study---21 cases and 9 lemmas, closed autonomously by two different models
for a few cents---demonstrates the workflow on exactly the kind of goal the
oracle must handle. The tool is open-source and under active development.

\textbf{Tool availability}: \url{https://github.com/scidonia/rocq-piler}

\bibliographystyle{splncs04}
\begin{thebibliography}{9}

\bibitem{coq}
{The Coq Development Team}:
The Coq Proof Assistant.
\url{https://doi.org/10.5281/zenodo.14542673} (2024)

\bibitem{jung2018iris}
Jung, R., Krebbers, R., Jourdan, J.-H., Bizjak, A., Birkedal, L., Dreyer, D.:
Iris from the ground up: A modular foundation for higher-order concurrent
separation logic.
Journal of Functional Programming \textbf{28}, e20 (2018)

\bibitem{first2023baldur}
First, E., Rabe, M.N., Ringer, T., Brun, Y.:
Baldur: Whole-Proof Generation and Repair with Large Language Models.
In: Proceedings of the 31st ACM Joint European Software Engineering Conference and Symposium on the Foundations of Software Engineering (ESEC/FSE) (2023)

\bibitem{yang2023leandojo}
Yang, K., Swope, A.M., Gu, A., Chalamala, R., Song, P., Yu, S., Godil, S., Prenger, R., Anandkumar, A.:
LeanDojo: Theorem Proving with Retrieval-Augmented Language Models.
In: Advances in Neural Information Processing Systems (NeurIPS) (2023)

\bibitem{gallego2022coqlsp}
Gallego Arias, E.J., et~al.:
Rocq-lsp: Visual Studio Code Extension and Language Server Protocol for Rocq, v0.2.5.
\url{https://github.com/ejgallego/rocq-lsp} (2025)

\bibitem{roe2016coqpie}
Roe, K.:
CoqPIE: A Plugin for Integrated Development Environments.
In: Conference on Intelligent Computer Mathematics (CICM) (2016)

\bibitem{aspinall2000proof}
Aspinall, D.:
Proof General: A Generic Tool for Proof Development.
In: Tools and Algorithms for the Construction and Analysis of Systems (TACAS) (2000)

\end{thebibliography}

\end{document}
